\chapter{Tiempo de Revisión y Análisis de Errores (OE3)}
\label{cap:revision}

% Ruta de las figuras y de las tablas que deja el cuaderno de este capítulo
% (notebook/figuras_06_oe3_revision.ipynb). Toda cifra \u... sale de ahí.
\newcommand{\figREV}{capitulos/06_oe3_revision/figures}

% ===================================================================
\section{Introducción}
% ===================================================================

Este capítulo presenta los resultados alcanzados para el tercer objetivo
específico:

\begin{quote}
    \textbf{OE3.} Cuantificar la reducción del tiempo de revisión que aporta el
    sistema y caracterizar sus errores mediante revisión experta de una muestra de
    predicciones.
\end{quote}

Los tres resultados comprometidos en la Tabla~\ref{tab:iov-oe3} ocupan una
sección cada uno: la metodología de validación comparativa (RE3.1,
sección~\ref{sec:re31}), la medición comparativa del tiempo de revisión (RE3.2,
sección~\ref{sec:re32}) y el análisis cualitativo de errores y hallazgos (RE3.3,
sección~\ref{sec:re33}). Cada sección sigue la estructura de los
Capítulos~\ref{cap:curacion} y~\ref{cap:deteccion}: el resultado, cómo se
alcanzó, cómo reproducirlo y la verificación de su indicador. La
sección~\ref{sec:discusion-oe3} discute los tres en conjunto.

El material de entrada es el modelo seleccionado en el
Capítulo~\ref{cap:deteccion}, con el umbral que allí se eligió en validación, y
la partición de prueba construida en el Capítulo~\ref{cap:curacion}. OE3 no
entrena ni vuelve a ajustar nada: mide lo que cuesta usar el sistema y de qué
están hechos sus errores. La medición la hicieron dos personas del equipo de
investigación entre el 11 y el 15 de noviembre de 2026, en dos tareas: una sesión
cronometrada de anotación manual y asistida (RE3.2) y la revisión ciega de
\uItems{} cajas del modelo (RE3.3).

% ===================================================================
\section{Metodología de validación comparativa (RE3.1)}
\label{sec:re31}
% ===================================================================

\rotulo{Resultado.} Un protocolo escrito que fija, antes de medir, qué se evalúa,
sobre qué datos, cuándo una predicción acierta, cómo se clasifica la que no
acierta y cómo se ejecutan las dos tareas de revisión experta: la sesión
cronometrada de RE3.2 y la revisión ciega de una muestra de cajas de RE3.3.

\rotulo{Cómo se alcanzó.} Con el procedimiento que la
Tabla~\ref{tab:metodos-resultados} asigna a este resultado: el protocolo se
redactó y se fijó en el repositorio, junto con el código que lo aplica, antes de
enviar el material al equipo. Su regla de acierto es la del
Capítulo~\ref{cap:deteccion}; lo nuevo es la unidad de análisis, la clasificación
de lo que no acierta y el diseño de las dos tareas.

\rotulo{Cómo reproducirlo.} El protocolo es \file{protocolo.md}, en la carpeta del
capítulo. El cuaderno \file{paquete\_06\_oe3\_revision.ipynb} arma, con la
semilla 42, las carpetas de la sesión y la muestra de cajas a partir de las tablas
que la herramienta produce sobre la partición de prueba, y
\file{src/analysis/review.py} contiene el emparejamiento y las categorías.

\rotulo{Medio de verificación.} El documento de metodología es el protocolo, y su
fecha, la del \textit{commit} que lo fijó en el repositorio. Esta sección lo
resume, y la sección~\ref{sec:verificacion-re31} contrasta el indicador elemento
por elemento.

\subsection{Qué se evalúa y sobre qué datos}

El sistema se evalúa tal como lo usaría el equipo de investigación: el modelo
\rWinner{} del Capítulo~\ref{cap:deteccion}, con el umbral de confianza de
\rThresholdYol{} elegido en validación, a través de la herramienta de inferencia
de RE2.4 (sección~\ref{sec:formato-salida}). La herramienta recorre cada
grabación en ventanas de 3~s con un salto de 1,5~s, funde las detecciones
repetidas entre ventanas solapadas y escribe una tabla de selección de Raven. La
unidad de análisis es una fila de esa tabla, es decir, una propuesta tal como la
ve el especialista, y no una detección por ventana como en el
Capítulo~\ref{cap:deteccion}; por eso las cifras de este capítulo no se comparan
una a una con las de aquel.

Los datos son las \uTestRecordings{} grabaciones de la partición de prueba
(\uTestHours{} horas de audio), que el modelo no vio ni al entrenar ni al elegir
el umbral. La referencia son las anotaciones del equipo limpiadas con el protocolo
de RE1.1, de las \nClasses{} clases del modelo y sin las filas que quedaron en
revisión. Las detecciones de las dos clases de ventana (\code{as/hc} y
\code{pt/dc}), cuya caja ocupa la ventana entera, quedan fuera. Sobre ese
material, la herramienta produce \uRows{} filas.

\subsection{Emparejamiento y categorías}

Una fila acierta si toma una anotación de su misma clase con un IoU de al menos
0,3, medido en segundos por escala mel, la misma escala de frecuencia que ve el
modelo. El emparejamiento es voraz: de mayor a menor confianza, cada fila toma la
anotación libre con que tiene mayor IoU, y cada anotación se toma una sola vez.
Es la regla del Capítulo~\ref{cap:deteccion} (sección~\ref{sec:protocolo-oe2})
aplicada a filas en lugar de ventanas.

Una fila sin acierto cae en la primera de cuatro reglas que cumpla, que
corresponden a la taxonomía de la Tabla~\ref{tab:taxonomia-fp}. Es \emph{caja
    distinta} si tiene un IoU de al menos 0,3 con una anotación de su clase que
otra fila ya tomó (tipo~B); \emph{otra llamada}, si lo tiene con una anotación de
otra clase (tipo~C); otra vez \emph{caja distinta}, si solapa con una anotación de
su clase con un IoU menor (tipo~A), y \emph{nada anotado}, si no cumple ninguna
(tipo~D). De las \uRows{} filas de prueba, \uMatchRows{} aciertan, \uBoxRows{} son
caja distinta, \uLabelRows{} otra llamada y \uNoneRows{} nada anotado. Otras
\uNotReviewed{} no tocan ninguna anotación de la referencia pero sí una que quedó
fuera de ella, en revisión o de una clase que el modelo no tiene; ahí ya hay algo
anotado, y no se revisan.

\subsection{Tarea 1: la sesión cronometrada}
\label{sec:diseno-sesion}

\rotulo{Material.} Cuatro carpetas de 14 grabaciones de prueba cada una, con dos
grabaciones de cada una de las siete especies que tienen algún tipo de llamada
entrenado; \textit{Cebus} queda fuera porque la condición asistida no le
propondría nada. Cada grabación dura a lo sumo 2~min y tiene al menos una
anotación. De 20\,000 sorteos se eligió el que dejaba las cuatro carpetas más
parecidas en minutos y en anotaciones: quedaron con entre \uFolderAudioLow{} y
\uFolderAudioHigh{} minutos y entre 47 y 51 anotaciones cada una, y suman
\uSessionMinutes{} minutos de audio.

\rotulo{Condiciones.} En la \emph{manual}, el especialista anota cada grabación
desde cero, como lo hace siempre. En la \emph{asistida}, parte de la tabla de la
herramienta, con la confianza de cada propuesta, y borra, corrige y agrega cajas
hasta dejarla como la entregaría. En las dos rigen el mismo estándar de calidad y
los mismos ajustes de Raven.

\rotulo{Orden.} Las carpetas se hacen en el orden manual, asistida, asistida,
manual, para que el aprendizaje a lo largo de la sesión recaiga parejo en las dos
condiciones. La sesión va antes que la revisión de cajas, cuando el especialista
todavía no conoce las propuestas del modelo, y usa grabaciones distintas.

\rotulo{Medidas.} El tiempo de trabajo de una carpeta es la hora de término menos
la de inicio, menos las pausas anotadas. Por condición se suman los tiempos y se
dividen por las horas de audio: minutos de trabajo por hora de audio, y el
cociente entre la condición asistida y la manual. No se fija ningún ahorro
esperado. Para que un tiempo menor no oculte un trabajo peor, cada tabla entregada
se compara con la referencia con la regla de acierto anterior: la
\emph{cobertura} es la parte de las anotaciones de referencia que alguna caja
entregada acierta, y la \emph{precisión}, la parte de las cajas entregadas que
aciertan alguna. En la condición asistida se registra además qué pasó con cada
propuesta, que se dejó igual (la caja entregada más parecida tiene un IoU de al
menos 0,9 y su misma clase), se corrigió (IoU de al menos 0,3) o se borró, y
cuántas cajas se agregaron (IoU menor que 0,3 con todas las propuestas). El costo
de la máquina se mide en minutos de cómputo por hora de audio, en GPU y en CPU.

\subsection{Tarea 2: la revisión ciega de cajas}
\label{sec:diseno-revision}

\rotulo{Muestra.} \uItems{} filas de prueba en \uItemsRecordings{} grabaciones,
sorteadas al azar dentro de seis estratos con la semilla 42. Las \uNoneRows{}
filas nada anotado se dividen en tres tramos de confianza (mayor que 0,5; de 0,3 a
0,5, y de 0,11 a 0,3), de los que se sortean 300, 200 y 133: el primero son las
predicciones de alta confianza que nombra el indicador, y los otros dos cubren el
resto del rango de operación. Las \uBoxRows{} filas de caja distinta y las
\uLabelRows{} de otra llamada entran completas, y \uControls{} de los
\uMatchRows{} aciertos entran como controles, para medir qué tan exigente es el
revisor con lo que la referencia ya da por bueno. Con 300 cajas, una proporción
tiene un intervalo de confianza del 95\,\% de $\pm$6 puntos en el peor caso, y
con las \uNoneSample{} de nada anotado juntas, de $\pm$4.

\rotulo{Ciego.} El revisor ve cada caja y su clase propuesta sobre el audio, en
una tabla de Raven que contiene solo las cajas de la muestra. No ve la confianza,
el estrato ni la referencia, y no sabe que hay controles.

\rotulo{Respuesta y lectura.} Cada caja recibe una sola palabra, más un comentario
opcional que dice cuál es la llamada o qué se oye. La Tabla~\ref{tab:lectura} da
el significado de cada palabra y cómo se lee según el estrato del que salió la
caja.

\begin{table}[htbp]
    \centering
    \small
    \begin{threeparttable}
        \caption{Respuestas del revisor y su lectura según el estrato}
        \label{tab:lectura}
        \begin{tabular}{@{}lL{4.1cm}L{2.5cm}L{3.0cm}L{2.6cm}@{}}
            \toprule
            \textbf{Respuesta} & \textbf{Significa} & \textbf{Nada anotado} & \textbf{Caja distinta u otra llamada} & \textbf{Acierto (control)} \\
            \midrule
            \code{correct} & Es la llamada propuesta y la caja la encuadra & Evento no anotado & Anotación por corregir & Acierto confirmado \\
            \addlinespace
            \code{box} & Es la llamada propuesta, con la caja mal puesta & Evento no anotado & Error de encuadre & Desacuerdo con la referencia \\
            \addlinespace
            \code{other} & Es una llamada de primate de otra especie o de otro tipo & Evento no anotado & Error de clase & Desacuerdo con la referencia \\
            \addlinespace
            \code{no} & No hay ninguna llamada de primate en la caja & Ruido & Ruido & Desacuerdo con la referencia \\
            \addlinespace
            \code{?} & No se puede decir & Incierto & Incierto & Incierto \\
            \bottomrule
        \end{tabular}
        \begin{tablenotes}[flushleft]
            \footnotesize
            \item \textit{Nota.} Si la clase y la caja están mal a la vez, la
            respuesta es \code{other}.
        \end{tablenotes}
    \end{threeparttable}
\end{table}

Una llamada real en una fila nada anotado es un evento que la referencia omitió,
aunque el modelo la haya encuadrado o etiquetado mal. En cambio, cuando el modelo
y la referencia discrepan sobre una llamada anotada, la respuesta dice cuál de los
dos tiene razón: si el revisor da la razón al modelo, lo que hay que corregir es la
anotación. De las categorías que no son acierto, tres son las que compromete el
indicador: evento no anotado, ruido y error de encuadre. El error de clase no cabe
en ninguna de ellas y se reporta aparte, como lo incierto, que no se reparte entre
las demás.

\rotulo{Estimación.} Por estrato, la proporción de cada respuesta con su intervalo
de Wilson al 95\,\% \citep{wilson_probable_1927}. Para las filas nada anotado en conjunto, cada tramo pesa lo
que pesa en la población (estimador estratificado, con intervalo normal), y de ahí
sale el número estimado de eventos no anotados en la partición de prueba y por
hora de audio. Si dos personas responden por separado las mismas cajas, su
acuerdo se mide con el kappa de Cohen; si no, se declara como limitación.

\subsection{Fijado antes de medir}

El protocolo y el código que lo aplica quedaron en el repositorio con el
\textit{commit} \code{63f84a8}, del 10 de octubre de 2026. Esa versión ya contiene
todas las reglas anteriores: el modelo y su umbral, la regla de acierto, las
categorías, los estratos y sus tamaños, el ciego, el formulario, la lectura, la
estimación y el diseño de la sesión, y su código ya aplica los umbrales de 0,9 y
0,3 con que se clasifica el destino de cada propuesta. El 27 de octubre, al armar
el material que se envió al equipo, el texto se reordenó y se fijó que la sesión
fuera la primera tarea. La sesión se hizo entre el 11 y el 13 de noviembre, y la
revisión de cajas terminó el 15. Ninguna regla de medición cambió después de ver
las respuestas.

\subsection{Desviaciones en la ejecución}
\label{sec:desviaciones-oe3}

La ejecución se apartó del protocolo en cuatro puntos. Ninguno cambia una regla,
pero los dos primeros cambian cómo se leen algunas cifras, y la
sección~\ref{sec:discusion-oe3} los discute.

\begin{enumerate}
    \item \textbf{Dos personas en lugar de una.} El protocolo preveía un solo
          revisor, y las dos tareas las hicieron dos personas del equipo, que aquí se
          llaman A y B. Según la nota que acompañó a los tiempos, B había clasificado
          antes vocalizaciones de las ocho especies, y A solo había trabajado con las
          del mono aullador; ninguna de las dos había anotado antes las grabaciones
          de la sesión. En la sesión se repartieron las carpetas de forma desigual
          (sección~\ref{sec:sesion}). En la revisión de cajas, cada caja tiene una
          sola respuesta, y las grabaciones volvieron separadas por persona en
          \uPersonRecordings{} de los \uItemsRecordings{} casos (\uPersonRecA{} de A
          y \uPersonRecB{} de B); en el resto no consta quién las revisó.
    \item \textbf{Revisión dentro de los intervalos.} Los tiempos de las carpetas
          1, 2 y 3 incluyen, cada uno, 15 minutos en que B revisó lo hecho por A.
    \item \textbf{Erratas en la planilla.} Las respuestas volvieron en una hoja de
          cálculo, no en el archivo de texto del paquete, y tres traían una errata:
          \code{corect} se leyó como \code{correct}; \code{other} con un espacio
          final, como \code{other}, y \code{too short}, escrita en la columna de la
          respuesta, como \code{box}, que es la respuesta a la que ese comentario
          acompaña en el resto de la planilla.
    \item \textbf{Un comentario fuera del formulario.} En todas las filas de
          \uRefrainRecordings{} de las \uItemsRecordings{} grabaciones, el
          comentario incluye «missing selection of vocalizations». Responde a una
          indicación que el autor dio por mensaje durante la tarea, fuera del
          protocolo: anotar también lo que el modelo omitiera. Los revisores no
          agregaron cajas, porque las tablas volvieron sin cambios, y dejaron esa
          nota. La tabla que veía
          el revisor tenía solo las cajas de la muestra, \uShownPerRecording{} por
          grabación en promedio, de las \uRowsPerSampledRecording{} filas que la
          herramienta produjo en esas grabaciones. El comentario describe esa tabla,
          en la que siempre quedaban llamadas sin caja, y no lo que el modelo omite;
          por eso no se usa.
\end{enumerate}

\subsection{Verificación de RE3.1}
\label{sec:verificacion-re31}

El indicador exige que el protocolo especifique seis elementos y que el punto de
operación y el umbral de IoU queden fijados antes de medir. La
Tabla~\ref{tab:verificacion-re31} localiza cada uno.

\begin{table}[htbp]
    \centering
    \small
    \begin{threeparttable}
        \caption{Verificación de RE3.1}
        \label{tab:verificacion-re31}
        \begin{tabular}{@{}L{4.3cm}L{3.0cm}L{6.0cm}@{}}
            \toprule
            \textbf{Elemento del IOV} & \textbf{Dónde} & \textbf{Qué fija el protocolo} \\
            \midrule
            Subconjunto de prueba & Sección~\ref{sec:re31} & \uTestRecordings{} grabaciones, \uTestHours{}~h, no vistas al entrenar ni al elegir el umbral \\
            Punto de operación elegido en validación & Sección~\ref{sec:re31} & Confianza de al menos \rThresholdYol{}, la del Capítulo~\ref{cap:deteccion} \\
            Umbral de IoU para verdadero positivo & Sección~\ref{sec:re31} & 0,3 con la misma clase \\
            Fijados antes de medir & Sección~\ref{sec:re31} & \textit{Commit} del 10-10-2026; sesión del 11 al 13-11-2026 \\
            Métricas & Secciones~\ref{sec:diseno-sesion} y~\ref{sec:diseno-revision} & Minutos por hora de audio y su cociente; cobertura y precisión; proporción por estrato con intervalo de Wilson; estimador estratificado \\
            Emparejamiento & Sección~\ref{sec:re31} & Voraz por confianza, cada anotación una sola vez, y cuatro reglas para lo que no acierta \\
            Procedimiento para revisar discrepancias y eventos potencialmente omitidos & Sección~\ref{sec:diseno-revision}, Tabla~\ref{tab:lectura} & Revisión ciega de \uItems{} cajas en seis estratos, con controles \\
            \bottomrule
        \end{tabular}
        \begin{tablenotes}[flushleft]
            \footnotesize
            \item \textit{Nota.} Protocolo: \file{protocolo.md}, con las reglas
            fijadas en el \textit{commit} \code{63f84a8}.
        \end{tablenotes}
    \end{threeparttable}
\end{table}

Cada elemento está en el protocolo con un valor concreto, y los dos que el
indicador exige fijar antes de medir lo estaban un mes antes de que el equipo
abriera la primera carpeta de la sesión.

% ===================================================================
\section{Medición comparativa del tiempo de revisión (RE3.2)}
\label{sec:re32}
% ===================================================================

\rotulo{Resultado.} El tiempo de revisión manual y asistida sobre el mismo
material, con la calidad de lo que se entregó en cada condición y el costo de
cómputo del modelo. En esta sesión, la revisión manual costó \uMinManual{}
minutos por hora de audio y la asistida \uMinAssisted{}, un cociente de
\uTimeRatio{}: la asistencia no redujo el tiempo. Las tablas asistidas, en
cambio, salieron más completas y más precisas.

\rotulo{Cómo se alcanzó.} Con la tarea~1 del protocolo
(sección~\ref{sec:diseno-sesion}), entre el 11 y el 13 de noviembre de 2026. El
equipo anotó en una planilla la hora de inicio, la de término y las pausas de cada
carpeta, y guardó en Raven una tabla por grabación.

\rotulo{Cómo reproducirlo.} El cuaderno \file{figuras\_06\_oe3\_revision.ipynb}
lee la planilla de tiempos y las tablas devueltas, las compara con la referencia y
genera las Tablas~\ref{tab:sesion} y~\ref{tab:calidad-sesion} y todas las cifras
de esta sección. Las tablas del modelo que recibió el equipo son las del paquete
de RE3.1.

\rotulo{Medio de verificación.} El informe de validación es esta sección. Los
registros de tiempo son la planilla del equipo, que la Tabla~\ref{tab:sesion}
resume carpeta por carpeta, y la medición de inferencia es la del
Capítulo~\ref{cap:deteccion} en GPU (sección~\ref{sec:costo-inferencia}) y la de
la herramienta en CPU sobre toda la partición de prueba. La
sección~\ref{sec:verificacion-re32} contrasta el indicador.

\subsection{La sesión}
\label{sec:sesion}

Las dos personas se repartieron las carpetas como indica su nota. En la
carpeta~1, manual, B anotó las seis primeras grabaciones y A las ocho restantes;
las carpetas~2 y~3, asistidas, las hizo sobre todo A, y en la carpeta~4, manual, A
anotó las cinco primeras y B las nueve restantes. La condición asistida quedó así
casi entera en manos de A, que tenía menos experiencia con estas especies, y la
manual, repartida entre las dos. De las \uSessionRecordings{} tablas devueltas,
\uEmptyTables{} llegaron vacías (\uEmptyManual{} manuales y \uEmptyAssisted{}
asistidas): según el equipo, en esas grabaciones no encontraron llamadas, aunque
cada una tiene al menos una anotación de referencia. Se cuentan tal como llegaron.

\subsection{Tiempo por condición}

La Tabla~\ref{tab:sesion} da el tiempo de trabajo por carpeta y condición.

\begin{table}[htbp]
    \centering
    \small
    \begin{threeparttable}
        \caption{Tiempo de trabajo en la sesión cronometrada}
        \label{tab:sesion}
        \begin{tabular}{@{}llrrrrrr@{}}
            \toprule
            \textbf{Carpeta} & \textbf{Condición} & \thead{Graba-\\ciones} & \thead{Audio\\(min)} & \thead{Anota-\\ciones} & \thead{Propues-\\tas} & \thead{Trabajo\\(min)} & \thead{Minutos por\\hora de audio} \\
            \midrule
            \input{\figREV/tablas/sesion} \\
            \bottomrule
        \end{tabular}
        \begin{tablenotes}[flushleft]
            \footnotesize
            \item \textit{Nota.} Carpetas en el orden en que se hicieron.
            Anotaciones: las de referencia. Trabajo: fin menos inicio, sin las
            pausas.
        \end{tablenotes}
    \end{threeparttable}
\end{table}

Las dos condiciones costaron casi lo mismo: \uWorkManual{} y \uWorkAssisted{}
minutos de trabajo sobre la misma cantidad de audio. Pero la diferencia entre
condiciones es pequeña al lado de la que hay entre las carpetas de una misma
condición: las dos manuales costaron \uFolderMinHigh{} y \uFolderMinLow{} minutos
por hora de audio, más de tres veces de diferencia, y el tiempo por hora bajó de
una carpeta a la siguiente, en el orden en que se hicieron. El orden manual,
asistida, asistida, manual reparte ese descenso entre las dos condiciones cuando
es parejo; con un descenso tan marcado, y con las personas repartidas de forma
distinta en cada carpeta, el cociente de \uTimeRatio{} no separa el efecto de la
asistencia del de quién anotó y en qué momento de la sesión.

\subsection{Calidad de lo entregado}

El mismo tiempo produjo tablas distintas. La Tabla~\ref{tab:calidad-sesion}
compara cada tabla entregada con la referencia, y agrega lo que habrían dado las
tablas del modelo sin tocar sobre las mismas grabaciones.

\begin{table}[htbp]
    \centering
    \small
    \begin{threeparttable}
        \caption{Calidad de las tablas de la sesión frente a la referencia}
        \label{tab:calidad-sesion}
        \begin{tabular}{@{}lcc@{}}
            \toprule
            \textbf{Tabla} & \textbf{Cobertura} & \textbf{Precisión} \\
            \midrule
            Entregada en la condición manual & \uRecallManual{} & \uPrecisionManual{} \\
            Entregada en la condición asistida & \uRecallAssisted{} & \uPrecisionAssisted{} \\
            \addlinespace
            Del modelo, en las grabaciones de la condición manual & \uModelRecallManual{} & \uModelPrecisionManual{} \\
            Del modelo, en las grabaciones de la condición asistida & \uModelRecallAssisted{} & \uModelPrecisionAssisted{} \\
            \bottomrule
        \end{tabular}
        \begin{tablenotes}[flushleft]
            \footnotesize
            \item \textit{Nota.} Referencia: \uSessionRef{} anotaciones de las
            clases del modelo en las \uSessionRecordings{} grabaciones. Acierto:
            misma clase e IoU de al menos 0,3.
        \end{tablenotes}
    \end{threeparttable}
\end{table}

Con las propuestas, las tablas cubrieron \uRecallAssisted{} de la referencia con
una precisión de \uPrecisionAssisted{}; sin ellas, \uRecallManual{} con
\uPrecisionManual{}. Las dos quedan lejos de la cobertura de las tablas del
modelo solo, que encuentran más anotaciones de referencia que cualquier tabla
humana a costa de una precisión menor. Las tablas humanas se miden contra una
referencia que otras personas hicieron en otro momento, lo que puede sesgar la
comparación \citep{koga_lead_2024}, y las \uEmptyTables{}
tablas vacías cuentan en contra de las dos condiciones.

Lo que se hizo con las propuestas explica parte de esa diferencia. De las
\uProposals{} propuestas que recibió la condición asistida, se dejaron igual
\uProposalsKeptN{} (\uProposalsKept{}), se corrigieron \uProposalsEditedN{}
(\uProposalsEdited{}) y se borraron \uProposalsDeletedN{} (\uProposalsDeleted{}).
De las \uProposalsHit{} que acertaban a la referencia, se dejaron igual
\uProposalsHitKept{}, se corrigieron \uProposalsHitEdited{} y se borraron
\uProposalsHitDeleted{}. Además se agregaron \uBoxesAdded{} cajas, de las que
\uBoxesAddedKnown{} tienen una clase del modelo y \uBoxesAddedHits{} aciertan a la
referencia. La revisión asistida, tal como se hizo, consistió sobre todo en
borrar: dos de cada tres propuestas, incluidas muchas correctas.

\subsection{Costo del modelo y lo que entrega}

La máquina cuesta poco al lado del especialista. La herramienta procesa una hora
de audio en \uCpuMinPerHour{} minutos de CPU, medidos sobre toda la partición de
prueba, y en \uGpuMinPerHour{} minutos de GPU (sección~\ref{sec:costo-inferencia}),
frente a los más de mil minutos por hora de audio que costó cualquiera de las dos
condiciones. Sobre toda la partición de prueba entrega \uProposalsPerHour{} filas
por hora de audio (\uProposalsPerHourBox{} sin las clases de ventana), y el
\uAudioWithoutProposals{} del tiempo de audio queda sin ninguna propuesta. La
condición asistida de la sesión recibió \uSessionProposalsPerHour{} propuestas por
hora de audio.

\subsection{Verificación de RE3.2}
\label{sec:verificacion-re32}

El indicador pide cuatro cosas y prohíbe una. La Tabla~\ref{tab:verificacion-re32}
las localiza.

\begin{table}[htbp]
    \centering
    \small
    \begin{threeparttable}
        \caption{Verificación de RE3.2}
        \label{tab:verificacion-re32}
        \begin{tabular}{@{}L{4.6cm}L{3.4cm}L{5.3cm}@{}}
            \toprule
            \textbf{Elemento del IOV} & \textbf{Dónde} & \textbf{Valor} \\
            \midrule
            Minutos por hora de audio, condición manual y asistida & Tabla~\ref{tab:sesion} & \uMinManual{} y \uMinAssisted{}; cociente \uTimeRatio{} \\
            Material comparable & Sección~\ref{sec:diseno-sesion}, Tabla~\ref{tab:sesion} & Cuatro carpetas de 14 grabaciones, de \uFolderAudioLow{} a \uFolderAudioHigh{} min y de 47 a 51 anotaciones cada una \\
            Orden de revisión & Secciones~\ref{sec:diseno-sesion} y~\ref{sec:sesion} & Manual, asistida, asistida, manual, antes de la revisión de cajas \\
            Tiempo de procesamiento del modelo por hora de audio & Sección~\ref{sec:re32} & \uCpuMinPerHour{} min en CPU y \uGpuMinPerHour{} min en GPU \\
            No fija un porcentaje de ahorro antes de medirlo & Sección~\ref{sec:diseno-sesion} & El protocolo no fija ningún ahorro esperado \\
            \bottomrule
        \end{tabular}
        \begin{tablenotes}[flushleft]
            \footnotesize
            \item \textit{Nota.} Tiempos de la planilla del equipo. Desviaciones:
            sección~\ref{sec:desviaciones-oe3}.
        \end{tablenotes}
    \end{threeparttable}
\end{table}

El indicador se cumple: comprometió medir el ahorro, no un valor para él, y la
medición se reporta con su material, su orden y el costo de la máquina.

% ===================================================================
\section{Análisis cualitativo de errores y hallazgos (RE3.3)}
\label{sec:re33}
% ===================================================================

\rotulo{Resultado.} La clasificación experta de \uItems{} filas del modelo en
\uItemsRecordings{} grabaciones de prueba, \uTopSample{} de ellas predicciones de
alta confianza sin anotación, en las categorías del indicador: \uCatUnannotated{}
eventos no anotados, \uCatNoise{} casos de ruido y \uCatFraming{} errores de
encuadre, más \uCatClass{} errores de clase, que van aparte. Llevada a toda la
partición de prueba, dice que tres de cada cuatro filas nada anotado son llamadas
reales que la referencia no tiene, y que solo una de cada cinco filas que la
referencia cuenta como error es ruido. La acompañan ejemplos en espectrograma de
cada categoría.

\rotulo{Cómo se alcanzó.} Con la tarea~2 del protocolo
(sección~\ref{sec:diseno-revision}). Las dos personas se repartieron las
grabaciones, respondieron la planilla de arriba abajo y la devolvieron completa el
15 de noviembre de 2026.

\rotulo{Cómo reproducirlo.} \file{paquete\_06\_oe3\_revision.ipynb} vuelve a
sortear la misma muestra con la semilla 42, y
\file{figuras\_06\_oe3\_revision.ipynb} lee la planilla devuelta, aplica la
lectura y la estimación del protocolo y genera la Tabla~\ref{tab:estratos}, la
Figura~\ref{fig:ejemplos} y todas las cifras de la sección.

\rotulo{Medio de verificación.} El informe de análisis es esta sección, y los
espectrogramas anotados, la Figura~\ref{fig:ejemplos}. La
sección~\ref{sec:verificacion-re33} contrasta el indicador.

\subsection{Respuestas por estrato}

La Tabla~\ref{tab:estratos} resume las respuestas por estrato.

\begin{table}[htbp]
    \centering
    \footnotesize
    \setlength{\tabcolsep}{3.5pt}
    \begin{threeparttable}
        \caption{Respuestas del revisor por estrato}
        \label{tab:estratos}
        \begin{tabular}{@{}L{3.6cm}rrccccc@{}}
            \toprule
            \textbf{Estrato} & \thead{Pobla-\\ción} & \thead{Revi-\\sadas} & \code{correct} & \code{box} & \code{other} & \code{no} & \thead{Llamada\\real} \\
            \midrule
            \input{\figREV/tablas/estratos} \\
            \bottomrule
        \end{tabular}
        \begin{tablenotes}[flushleft]
            \footnotesize
            \item \textit{Nota.} Porcentaje de las revisadas y, entre paréntesis,
            intervalo de Wilson al 95\,\%. Los tramos de nada anotado son de
            confianza; su fila conjunta pondera por población (intervalo normal).
            Llamada real: \code{correct}, \code{box} u \code{other}. Las
            \uUnsure{} respuestas \code{?} no figuran.
        \end{tablenotes}
    \end{threeparttable}
\end{table}

\rotulo{Nada anotado.} En el tramo de alta confianza, \uNoneRealTop{} de las
cajas (\uNoneRealTopLow{} a \uNoneRealTopHigh{}) encuadran una llamada real, y
\uNoneCorrectTop{} son exactamente la llamada propuesta con su caja: son
detecciones correctas que la referencia no tiene. La parte de llamadas reales baja
con la confianza, de \uRealLastBin{} por encima de 0,7 a \uRealFirstBin{} entre
0,11 y 0,2, y el ruido sube de \uNoiseLastBin{} a \uNoiseFirstBin{}: la confianza
ordena las propuestas, aunque no separa las reales del ruido. Para las
\uNoneRows{} filas nada anotado en conjunto, el estimador estratificado da
\uNoneReal{} de llamadas reales (de \uNoneRealLow{} a \uNoneRealHigh{}). Son unas
\uUnannotated{} filas de la partición de prueba (de \uUnannotatedLow{} a
\uUnannotatedHigh{}), \uUnannotatedPerHour{} por hora de audio, sobre
vocalizaciones que nadie anotó, frente a las \uRefTest{} anotaciones de la
referencia en las mismas clases, \uRefPerHour{} por hora. El ruido es el
\uNoneNoise{} restante (de \uNoneNoiseLow{} a \uNoneNoiseHigh{}), unas
\uNoisePerHour{} filas por hora de audio.

\rotulo{Caja distinta y otra llamada.} Donde el modelo y la referencia marcan la
misma clase sin coincidir en la caja, el revisor juzgó mal puesta la caja del
modelo en \uBoxModelFraming{} de los casos y bien puesta en \uBoxModelOk{}; en
estos últimos, la que hay que corregir es la anotación. Donde el modelo propone
otra clase sobre una llamada anotada, el revisor dio la razón a la referencia en
\uLabelModelWrong{} de los casos y al modelo en \uLabelModelOk{}. La confusión
depende de entre qué clases ocurre: \uLabelSameSpecies{} de esas \uLabelRows{}
filas tocan una anotación de la misma especie, es decir, de otro tipo de llamada,
y en las \uLabelCross{} restantes, que cruzan especies, el revisor confirmó que la
clase del modelo era otra en \uLabelCrossOther{}.

\rotulo{Controles.} De los \uControls{} aciertos, el revisor confirmó
\uControlsOk{}. En \uControlsFraming{} juzgó mal puesta la caja, aunque coincide
con la anotación en un IoU de al menos 0,3, y en \uControlsNoiseN{} casos no oyó
ninguna llamada de primate (\uControlsBird{} de ellos los describió como ave). El
revisor es, entonces, más exigente con el encuadre que la regla de acierto, y la
propia referencia tiene errores.

\subsection{Qué es el error aparente}

La precisión que se mide contra la referencia trata como error toda fila que no
acierta. Si cada estrato se lleva a su población, las \uApparent{} filas
revisables que no aciertan se reparten así según el revisor: en \uApparentCorrect{}
el modelo tenía razón; en \uApparentBox{} había una llamada real con la caja mal
puesta; en \uApparentOther{}, una llamada real de otra clase, y solo en
\uApparentNoise{} no había ninguna llamada de primate. Con el revisor como
referencia, la precisión por fila, que contra la referencia es \uPrecisionRef{},
sube a \uPrecisionStrict{} si se exigen la clase y la caja correctas, a
\uPrecisionLenient{} si basta que la llamada de esa clase esté en la caja y a
\uPrecisionCall{} si basta cualquier llamada de primate. Esta lectura no es una
regla del protocolo: combina sus proporciones por estrato con el tamaño de cada
estrato, aplica a los \uMatchRows{} aciertos lo que el revisor respondió en los
controles y deja fuera las \uNotReviewed{} filas que no se revisaron.

\subsection{Errores comunes}

\rotulo{El ruido.} De las \uNoAnswers{} respuestas \code{no}, el comentario
describe un ave en \uNoBird{}, ruido de fondo en \uNoBackground{}, insectos en
\uNoInsect{} y voces o ruidos humanos en \uNoHuman{}; el resto no lo describe. El
ruido se concentra en pocas clases: \uNoiseTopClass{} aporta \uNoiseTopK{} de sus
\uNoiseTopN{} cajas nada anotado revisadas. Por especie, la parte de ruido entre
las cajas nada anotado es \uNoiseBySpecies{}.

\rotulo{Las confusiones de clase.} Los comentarios de las respuestas \code{other}
nombran sobre todo otro tipo de llamada de la misma especie: los tipos de SB entre
sí, los de LW con su \code{trino} y, en AS, \code{bc} con \code{hc}. La clase que
más se confunde es \uOtherTopClass{}: \uOtherTopK{} de sus \uOtherTopN{} cajas
nada anotado son, según el revisor, otra llamada. Es lo que anticipaba el
Capítulo~\ref{cap:deteccion}: la clase de especie se distingue bien, y el error
está en el tipo de llamada y en la caja.

\subsection{Ejemplos}

La Figura~\ref{fig:ejemplos} muestra un ejemplo de cada categoría, elegido con
una regla fija: la caja de mayor confianza de la categoría y, en el ruido, la de
mayor confianza entre las que el revisor describe como ave, la fuente de ruido
más frecuente.

\begin{figure}[htbp]
    \centering
    \caption{Ejemplos de la revisión de cajas}
    \label{fig:ejemplos}
    \includegraphics[width=\textwidth]{\figREV/revision_ejemplos.png}
    \par\medskip
    \begin{minipage}{\textwidth}
        \footnotesize
        \textit{Nota.} Ventanas de 3~s centradas en la caja, en escala mel. Entre
        paréntesis, la confianza. «Anotada»: clase de la anotación que toca la caja,
        si difiere; «revisor»: la llamada que nombró el revisor.
    \end{minipage}
\end{figure}

El acierto confirmado y la llamada no anotada se parecen: en los dos, una caja
ajustada sobre una llamada nítida, pero en el segundo no hay anotación y la
referencia lo cuenta como falso positivo.

\subsection{Sensibilidad}

Dos lecturas adicionales, que no cambian ninguna regla, acotan la cifra principal.
La primera se refiere a las cajas nada anotado que caen dentro de una anotación de
otra clase, por lo general más larga, como un coro: \uNoneInside{} de las
revisadas tienen al menos la mitad de su área dentro de una sola anotación, y
\uNoneInsideOther{} de ellas el revisor las respondió \code{other}, casi siempre
nombrando la clase de esa anotación. La lectura del protocolo las cuenta como
eventos no anotados; sin ellas, la parte de llamadas reales baja de \uNoneReal{} a
\uNoneRealKept{}, y los eventos no anotados por hora de audio, de
\uUnannotatedPerHour{} a \uUnannotatedKeptPerHour{}.

La segunda se refiere a los intervalos, que suponen cajas independientes, cuando
las de una misma grabación se parecen. Remuestreando grabaciones dentro de cada
estrato, el intervalo del estimador estratificado pasa a ir de
\uNoneRealBootLow{} a \uNoneRealBootHigh{}, apenas más ancho. El de la respuesta
\code{other} en otra llamada, en cambio, pasa de ir de \uLabelOtherLow{} a
\uLabelOtherHigh{} a ir de \uLabelOtherBootLow{} a \uLabelOtherBootHigh{}, porque
\uLabelLargestRecording{} de sus \uLabelRows{} filas vienen de una sola grabación.
Las cifras de nada anotado, que son las que responden a la causa C3, resisten las
dos lecturas.

\subsection{Verificación de RE3.3}
\label{sec:verificacion-re33}

El indicador pide clasificar una muestra de predicciones de alta confianza en tres
categorías y mostrar tres clases de ejemplos. La Tabla~\ref{tab:verificacion-re33}
localiza cada uno.

\begin{table}[htbp]
    \centering
    \small
    \begin{threeparttable}
        \caption{Verificación de RE3.3}
        \label{tab:verificacion-re33}
        \begin{tabular}{@{}L{4.5cm}L{3.3cm}L{5.5cm}@{}}
            \toprule
            \textbf{Elemento del IOV} & \textbf{Dónde} & \textbf{Valor} \\
            \midrule
            Muestra de predicciones de alta confianza & Tabla~\ref{tab:estratos} & \uTopSample{} con confianza mayor que 0,5, dentro de \uItems{} revisadas \\
            Evento no anotado & Tabla~\ref{tab:estratos} & \uCatUnannotated{} en la muestra; \uNoneReal{} de las filas nada anotado de prueba \\
            Ruido & Tabla~\ref{tab:estratos} & \uCatNoise{} en la muestra; \uNoneNoise{} de las filas nada anotado de prueba \\
            Error de encuadre & Tabla~\ref{tab:estratos} & \uCatFraming{} en la muestra; \uBoxModelFraming{} de las filas caja distinta \\
            Ejemplos de aciertos & Figura~\ref{fig:ejemplos} & Acierto confirmado \\
            Ejemplos de errores comunes & Figura~\ref{fig:ejemplos} & Ruido, error de encuadre y error de clase \\
            Ejemplos de detecciones correctas ausentes de la referencia & Figura~\ref{fig:ejemplos} & Llamada no anotada \\
            \bottomrule
        \end{tabular}
        \begin{tablenotes}[flushleft]
            \footnotesize
            \item \textit{Nota.} La muestra tiene además \uCatClass{} errores de
            clase, \uCatFix{} anotaciones por corregir, \uCatConfirmed{} aciertos
            confirmados, \uCatDisagree{} desacuerdos en los controles y
            \uCatUnsure{} respuestas inciertas.
        \end{tablenotes}
    \end{threeparttable}
\end{table}

La muestra cubre el tramo de alta confianza que nombra el indicador y, además, el
resto del rango de operación, y cada una de las \uItems{} cajas quedó en una
categoría. El indicador se cumple.

% ===================================================================
\section{Discusión}
\label{sec:discusion-oe3}
% ===================================================================

\rotulo{Resumen.} Las reglas de la evaluación quedaron fijadas un mes antes de
medir (sección~\ref{sec:re31}). En la sesión cronometrada, revisar con las
propuestas del modelo tomó lo mismo que anotar desde cero, pero dio tablas más
completas y más precisas (sección~\ref{sec:re32}). Y la revisión de cajas mostró
que la mayor parte de lo que la referencia cuenta como error del modelo no lo es:
las predicciones sin anotación son, en su mayoría, llamadas que nadie anotó, y el
ruido es una parte menor y concentrada en pocas fuentes (sección~\ref{sec:re33}).

\rotulo{Qué significan.} Para el objetivo general, esta sesión no muestra que el
sistema reduzca el tiempo de revisión. Muestra, en cambio, en qué se fue el tiempo
asistido: en borrar dos de cada tres propuestas, incluidas \uProposalsHitDeleted{}
de las \uProposalsHit{} que acertaban, y en agregar cajas nuevas. Con la revisión
de cajas al lado, ese borrado no se explica por la calidad de las propuestas, que
en su mayoría encuadran llamadas reales, sino por el criterio con que se juzga una
caja: incluso entre los aciertos, el revisor marcó \uControlsFraming{} como mal
encuadrados. Si una propuesta que no se ajusta al criterio del especialista se
borra y se vuelve a dibujar, la asistencia ahorra la búsqueda pero no el dibujo, y
en un material tan denso como el de la sesión, donde casi toda grabación tiene algo
que anotar, la búsqueda pesa poco. Es una explicación compatible con las cifras, no
una medición: la sesión no registró cuánto tiempo fue a cada acción.

Para la causa C3, el resultado es más claro. Medida contra la referencia, la
precisión del modelo por fila es \uPrecisionRef{}; según el revisor, está entre
\uPrecisionStrict{} y \uPrecisionLenient{}. La diferencia es la referencia: le
faltan unas \uUnannotatedPerHour{} llamadas por hora de audio que el modelo sí
propone, y algunas de las que tiene están mal encuadradas o mal etiquetadas. Eso
tiene dos consecuencias para el equipo. La primera es que las métricas del
Capítulo~\ref{cap:deteccion}, medidas contra esa misma referencia, subestiman al
modelo. La segunda es que no hace falta revisarlo todo por igual: las propuestas
de alta confianza sin anotación son llamadas reales en \uNoneRealTop{} de los
casos, y el ruido viene sobre todo de aves y de unas pocas clases, de modo que el
especialista puede decidir qué aceptar, qué revisar primero y qué filtrar, que es
lo que la causa C3 echaba en falta (sección~\ref{sec:causas}).

\rotulo{Frente a la literatura.} Los ahorros revisados en la
sección~\ref{sec:magnitud} corresponden a detectar una o pocas especies: depurar
los falsos positivos de un sistema semiautomático para primates costó entre el 3 y
el 5\,\% del tiempo del procesamiento humano \citep{heinicke_2015}, HOWLish redujo
15 veces el tiempo del operador \citep{Campos_Krofel_Rio-Maior_Renna_2025} y un
detector de llamadas de langur redujo en un 87\,\% la carga de revisión
\citep{hong_deep_2026}. Ninguno de esos trabajos medía la anotación completa, con
caja, especie y tipo de llamada, que es la que esta sesión cronometró
(sección~\ref{sec:vacio}). Que aquí no aparezca un ahorro es coherente con que
delimitar cada evento sea muy costoso en tiempo \citep{koga_lead_2024}: el detector resuelve la búsqueda, pero no el encuadre, y es
en el encuadre donde el especialista y el modelo discrepan. Esa discrepancia
también es la que anticipa la literatura. La anotación manual tiene errores
propios de cada anotador \citep{nguyen_assessing_2021}, dos analistas pueden
coincidir en menos de la mitad de los casos \citep{leroy_reliability_2018}, y la
borrosidad de las fronteras desplaza las cajas aun cuando el evento es el mismo
\citep{zhu_sound_2025}. El \uControlsFraming{} de aciertos que el revisor juzgó mal
encuadrados, y la distinta frecuencia con que las dos personas respondieron
\code{box}, son ese mismo fenómeno medido en este corpus. Por último, que un
detector sea útil con una precisión baja frente a la referencia tiene precedente:
HOWLish lo es con una precisión de 0,006 \citep{Campos_Krofel_Rio-Maior_Renna_2025}.

\rotulo{Generalización.} Lo que se sostiene fuera de este corpus es el método: una
muestra estratificada, revisada a ciegas y con controles, separa el error del
detector del de la referencia con cualquier detector y cualquier referencia, y el
estimador estratificado lleva la muestra a la población. Es probable que se
sostenga también la conclusión cualitativa de que una referencia hecha a mano
omite llamadas, porque la literatura lo documenta en otros dominios. Las cifras,
en cambio, no se sostienen sin volver a medir: los minutos por hora de audio
dependen de quién anota, de su experiencia y de la densidad del material, y aquí
los midieron dos personas sobre \uSessionMinutes{} minutos de audio elegido entre
grabaciones con anotaciones; las fuentes de ruido dependen del sitio y de la
temporada, y las confusiones de clase, del repertorio de cada especie.

\rotulo{Limitaciones.} Cada una dice cómo leer un resultado, no lo invalida.

\begin{itemize}
    \item \textbf{Dos personas con distinta experiencia.} La condición asistida la
          hizo casi entera A, la persona con menos experiencia en estas especies, y
          la manual se repartió entre las dos. El cociente de \uTimeRatio{} mezcla,
          por eso, el efecto de la asistencia con el de la persona, y los 15
          minutos de revisión dentro de tres intervalos suman tiempo a las dos
          condiciones sin distinguirlas.
    \item \textbf{Sin medida de acuerdo.} Cada caja tiene una sola respuesta, así
          que no hay kappa. En las \uPersonRecordings{} grabaciones que volvieron
          separadas por persona, A respondió \code{box} en \uPersonBoxA{} de las
          cajas que no eran controles y B en \uPersonBoxB{}; dentro de una misma
          especie, \uPersonSpecies{}, en \uPersonBoxSpeciesA{} y
          \uPersonBoxSpeciesB{}. La división entre \code{correct} y \code{box}
          depende, entonces, de quién revisó, y la cifra robusta es la que las suma:
          llamada real frente a ruido.
    \item \textbf{Una sesión corta y densa.} \uSessionMinutes{} minutos de audio en
          \uSessionRecordings{} grabaciones, todas con alguna anotación, y una
          variación entre carpetas mayor que la diferencia entre condiciones. El
          tiempo por hora de audio no se extrapola a una grabación continua ni a
          una temporada.
    \item \textbf{Filas que no son llamadas distintas.} Una llamada con dos filas
          cuenta dos veces, y el revisor, que veía solo las cajas de la muestra, no
          podía notar un duplicado. Los eventos no anotados por hora son filas, no
          llamadas.
    \item \textbf{Una referencia con errores.} Los controles muestran que la
          referencia también se equivoca; las cifras de calidad de la sesión, que
          se miden contra ella, heredan esos errores.
    \item \textbf{Lo que quedó fuera.} Las \uNotReviewed{} filas que tocan
          anotaciones fuera de la referencia no se revisaron, y la tarea opcional
          sobre la calidad del corpus entero, que no forma parte del indicador,
          queda para después de esta entrega.
    \item \textbf{Un punto de operación fijo.} El umbral de \rThresholdYol{} no se
          movió, por diseño. Uno más alto reduciría el ruido, que baja con la
          confianza, pero también las llamadas reales que el sistema propone.
\end{itemize}
